% !TEX root = ../../main.tex
\section{Giới thiệu đề tài}
\label{sec:gioi_thieu}
Đề tài tốt nghiệp \textit{``Xây dựng nền tảng AI hỗ trợ nghiên cứu và viết bài báo khoa học VinUni dựa trên Agentic RAG, truy xuất tri thức học thuật và kiểm chứng trích dẫn''} được nghiên cứu và hiện thực hóa thông qua hệ thống phần mềm \VS.

Mục tiêu trọng tâm của đề tài là xây dựng một nền tảng trợ lý trí tuệ nhân tạo hỗ trợ quy trình nghiên cứu khoa học lấy môi trường đại học định hướng nghiên cứu quốc tế như VinUniversity làm bối cảnh đối chuẩn điển hình, đồng thời phục vụ rộng rãi đội ngũ giảng viên, nghiên cứu sinh, học viên cao học và sinh viên đại học. Hệ thống liên kết xuyên suốt bốn không gian làm việc trong vòng đời của một công trình khoa học: từ khâu tìm kiếm khám phá ấn phẩm đa nguồn, phân tích đọc sâu tài liệu chuyên ngành hai cột, tổng hợp mạng lưới tri thức học thuật (GraphRAG \cite{edge2024}), cho đến khâu soạn thảo và biên dịch bản thảo chuẩn quốc tế bằng LaTeX.

Hệ thống \VS tập trung giải quyết bài toán giảm thiểu hiện tượng ảo giác trích dẫn trong nghiên cứu học thuật bằng cách kết hợp kiến trúc tác nhân tạo sinh tăng cường truy xuất (Agentic RAG -- \textit{Retrieval-Augmented Generation}) với hai tầng bảo chứng bổ trợ lẫn nhau: thứ nhất là đối chiếu tọa độ trực quan đoạn văn nguồn (\textit{visual citation grounding}) thông qua việc ánh xạ từng thẻ trích dẫn về đúng hộp tọa độ chuẩn hóa trên trang PDF gốc để học giả trực tiếp thẩm định bằng mắt thường; thứ hai là hậu kiểm tính bám sát bằng chứng (\textit{semantic and lexical claim verification}) nhằm từ chối sinh câu trả lời khi ngữ cảnh thu hồi không đủ độ tin cậy, qua đó hạn chế rủi ro dẫn chiếu sai lệch trong các ấn phẩm khoa học.
